% !TeX root = ../main.tex
\chapter{Methods}

\section{Study design}
Describe the study population, sampling, imaging procedures, outcomes and
analysis plan. Table~\ref{tab:measures} demonstrates an ordinary floating table.
LaTeX supplies the chapter-based number and adds the caption to the list of
tables.

For illustration, assume a prospective cohort with baseline and 24-month
follow-up. Participants would undergo symptom assessment, clinical examination
and a standardised imaging protocol. An analysis plan written before outcome
assessment would identify the primary measure, expected confounders and approach
to missing observations.

\subsection{Eligibility}
An example set of inclusion criteria might require adults with knee symptoms
but no established radiographic disease. Exclusion criteria could include
contraindications to MRI, inflammatory arthritis, recent major trauma or
inability to complete follow-up.

\subsubsection{Recruitment pathway}
Potential participants could be identified through primary care, community
advertising or existing research cohorts. Screening should use the same written
criteria in every setting, and the number approached, assessed and enrolled
should be retained for the participant-flow diagram.

\subsubsection{Consent and withdrawal}
The information sheet should explain the imaging procedure, use of clinical
data, foreseeable burdens and the right to withdraw. The thesis should
distinguish withdrawal from further contact, withdrawal from future data
collection and any request concerning data already obtained.

\subsection{Imaging acquisition}
Scanner field strength, coils, positioning, sequences, voxel dimensions and
quality-control procedures should be reported sufficiently clearly for the
method to be reproduced.

\begin{table}[htbp]
  \caption{Example research measures}
  \label{tab:measures}
  \centering
  \begin{tabular}{lll}
    \toprule
    Domain & Measure & Purpose \\
    \midrule
    Imaging & Cartilage thickness & Structural change \\
    Clinical & Pain score & Patient experience \\
    Mechanical & Knee load & Physical exposure \\
    \bottomrule
  \end{tabular}
\end{table}

\section{Multi-page tables}
The \texttt{longtable} environment repeats its headings automatically when a
table crosses a page boundary. The short example below uses the same structure
and can be extended with additional rows.

\begin{longtable}{@{}p{0.22\textwidth}p{0.28\textwidth}p{0.38\textwidth}@{}}
\caption{Example data dictionary}
\label{tab:data-dictionary}\\
\toprule
Variable & Type & Description \\
\midrule
\endfirsthead

\multicolumn{3}{@{}l}{\tablename\ \thetable\ continued}\\
\toprule
Variable & Type & Description \\
\midrule
\endhead

\midrule
\multicolumn{3}{r@{}}{Continued on next page}\\
\endfoot

\bottomrule
\endlastfoot

Participant ID & Text & Pseudonymous study identifier. \\
Age & Number & Age at baseline assessment. \\
Pain score & Number & Patient-reported symptom measure. \\
Cartilage thickness & Number & Quantitative structural imaging measure. \\
T2 relaxation time & Number & Quantitative compositional imaging measure. \\
\end{longtable}

\section{Analysis}
Continuous variables could be summarised using means and standard deviations or
medians and interquartile ranges, depending on their distributions. Model
assumptions should be checked and uncertainty reported using confidence
intervals rather than relying only on thresholded significance tests.

\begin{equation}
  Y_{24} = \beta_0 + \beta_1 X_{\mathrm{MRI}}
           + \beta_2 Y_0 + \beta_3 A + \varepsilon,
  \label{eq:analysis-model}
\end{equation}

Equation~\ref{eq:analysis-model} illustrates a follow-up outcome adjusted for
baseline outcome \(Y_0\), age \(A\), and an MRI predictor.

\subsection{Sensitivity analyses}
Sensitivity analyses might repeat the primary model after alternative handling
of missing data, exclusion of scans that failed quality control and adjustment
for additional prespecified confounders. These analyses should test the
robustness of interpretation rather than act as a search for a preferred
result.

\subsection{Reporting}
Results should be presented with denominators, effect estimates and confidence
intervals. Tables should be understandable without repeatedly consulting the
main text, while the narrative should highlight the pattern of findings rather
than reproduce every number in the table.
